% ==============================================================================
% Section: Capped Auctions — Quantity Rationing in 2026/27 and 2027/28
% Status: NEW (2026-05-02) — extracted from sec8_21billion.tex case-study
%   subsections; reframed around quantity-rationing channel; $21B
%   reconciliation moved to conclusion.
% ==============================================================================

\section{Capped Auctions: Quantity Rationing in 2026/27 and 2027/28}
\label{sec:case_studies}

The cap-incidence panel of Section~\ref{sec:cap_incidence} establishes
that a quantity shortfall against the reliability requirement emerged
in the two delivery years for which the Shapiro settlement cap was
binding. This section turns to the dollar and physical consequences of
the shortfall in each auction. Both cases use the IMM's mechanical
unrestricted-VRR counterfactual---which re-clears the same offers
against the pre-settlement demand curve---as the principal benchmark.

\label{subsec:case_2026}%
The 2026/27 BRA cleared at \$329.17/MW-day, which under the new VRR
design embodies the settlement cap. The market cleared 314~MW below
the RTO reliability requirement of 134{,}519.5~MW (a 0.23\% shortfall)
and 492~MW below the IRM-implied reserve target; the IMM declared the
auction ``not competitive'' \citep{MonitoringAnalytics2025_sotm}.
Capacity Performance offers had unit-specific Market Seller Offer Caps
calculated for 6.3\% of generation resources.

The IMM's mechanical counterfactual scenario re-clears the same offers
against the unrestricted (pre-settlement) VRR demand curve and yields
total revenues of \$19.29~billion, a 19.7\% increase over the
\$16.12~billion actual revenue \citep{MonitoringAnalytics2025_sotm}.
The system-level revenue increase substantially exceeds what the
\$4.17/MW-day RTO-level binding gap alone implies; most of the
difference comes from constrained Locational Deliverability Areas in
which local Point~(a) caps lie further above the settlement cap than
at the RTO level, so re-clearing those LDAs against the unrestricted
VRR raises both their local prices and the imports they pay for. With
the RTO reserve margin only marginally negative, the cap truncated
the price signal more than it truncated the quantity procured at the
system level.

\label{subsec:case_2027}%
The 2027/28 BRA cleared at \$333.44/MW-day, again at the new-design
Point~(a). The market cleared 6{,}623~MW below the RTO reliability
requirement of 141{,}101~MW---a 4.7\% shortfall that the IMM
characterized as material; the auction was declared ``not competitive''
on the same basis as 2026/27. Unit-specific MSOCs were calculated for
only 1.0\% of resources, consistent with the settlement cap moving
close enough to the underlying offer distribution that fewer individual
offers exceeded their cost-basis caps.

The IMM unrestricted-VRR counterfactual produces total revenues of
\$26.32~billion, a 60.4\% increase over the \$16.41~billion actual
revenue. The disproportionately large increase relative to 2026/27
reflects two compounding channels: a higher counterfactual price (the
unrestricted VRR allows clearing further up the demand curve) and a
higher counterfactual cleared quantity (offers above the settlement
cap that were excluded under the actual auction would have cleared).
A back-of-envelope decomposition: at a counterfactual price of
\$500/MW-day and the actual cleared quantity of 134{,}478~MW over 366
days, revenues would be approximately \$24.6~billion; the IMM's
\$26.32~billion exceeds this by roughly \$1.7~billion, implying a
counterfactual cleared quantity of approximately 144{,}000~MW---about
9{,}500~MW above actual and roughly equal to the reliability
requirement plus a positive reserve margin. Under the IMM scenario the
unrestricted auction would have realized higher prices and procured
adequate capacity; under the cap, the auction achieved neither.

\begin{table}[H]
  \centering
  \caption{Capacity Market Revenue Under Counterfactual Cap Regimes
    (System Total)}
  \label{tab:capped_revenue}
  \small
  \begin{tabular}{lcccc}
    \toprule
                                                 & 2026/27 & 2027/28 & Two-year & Implied \\
    Scenario                                     & revenue & revenue & total    & savings \\
    \midrule
    Actual (settlement-modified VRR)             & \$16.12B & \$16.41B & \$32.53B & --- \\
    \addlinespace
    IMM unrestricted-VRR counterfactual          & \$19.29B & \$26.32B & \$45.62B & \$13.08B \\
    \quad \textit{Percent increase vs.\ actual}  & +19.7\%  & +60.4\%  & +40.2\%  & --- \\
    \addlinespace
    \$500/MW-day at actual cleared MW            & \$24.51B & \$24.60B & \$49.10B & \$16.57B \\
    \addlinespace
    Stated Pennsylvania claim                    & ---      & ---      & $\sim$\$21B & $\sim$\$21B \\
    \bottomrule
  \end{tabular}
  \vspace{3pt}

  \begin{minipage}{0.96\textwidth}
    \footnotesize \textit{Notes:} Revenues in nominal dollars over the
    twelve-month delivery year (366 days for 2027/28). The IMM
    unrestricted-VRR counterfactual holds offers fixed and re-clears
    against the pre-settlement VRR demand curve; figures from the IMM
    2025 State of the Market Report
    \citep{MonitoringAnalytics2025_sotm}. The \$500/MW-day scenario
    multiplies the actual cleared UCAP (134{,}205~MW in 2026/27,
    134{,}478~MW in 2027/28) by Pennsylvania's stated counterfactual
    price; this scenario corresponds to the Pennsylvania counterfactual
    but holds quantity at the capped level rather than re-clearing
    demand. The \$21~billion claim
    \citep{Shapiro2025_press_release} reflects additional revenue
    contributions from constrained Locational Deliverability Areas
    where local Point~(a) prices exceeded \$325/MW-day in prior years.
    Implied savings are differences between the counterfactual scenario
    and the actual outcome.
  \end{minipage}
\end{table}

The 2027/28 case is the cleanest evidence in the sample of the
quantity-rationing channel of a binding administrative cap. The
2026/27 case isolates the price channel in a market only marginally
short of adequacy. Together the two cases trace the empirical
realization of the textbook prediction of
Section~\ref{subsec:theory_ceiling}: a binding ceiling truncates the
price signal in any tight market, and where supply at the cap is
insufficient to meet the planning target, it also truncates the
quantity procured.
